\chapter{Lesson 10: Irregular Verbs and Course Review}

Irregular verbs were introduced one at a time over the term (\ko{덥다} in the weather
lesson, \ko{듣다} with hobbies, \ko{아프다} at the hospital). In the last week we put them
all together. The good news is that Korean ``irregulars'' are really \emph{regular
classes}: every irregular verb belongs to one of a handful of patterns, defined by the
last letter of the stem, and the change only happens before endings that start with a
\textbf{vowel} (\ko{\til 아/어요}, \ko{\til 았/었어요}, \ko{\til (으)면},
\ko{\til (으)세요}, \ko{\til (으)ㄹ 거예요}). Before consonant-initial endings
(\ko{\til 고}, \ko{\til 지만}, \ko{\til 지 마세요}, \ko{\til 네요}) almost all of them
behave normally.

\section{How to tell which pattern applies}

\begin{center}
\begin{tikzpicture}[every node/.style={font=\small},
    level distance=1.7cm,
    level 1/.style={sibling distance=2.25cm},
    t/.style={draw, rounded corners, fill=blue!6, align=center, text width=1.9cm, font=\scriptsize}]
  \node[draw, rounded corners, fill=gray!15] {What does the stem end in?}
    child {node[t] {\ko{ㅂ}\\ \ko{ㅂ} $\to$ \ko{우}\\ \ko{더워요}}}
    child {node[t] {\ko{ㄷ}\\ \ko{ㄷ} $\to$ \ko{ㄹ}\\ \ko{들어요}}}
    child {node[t] {\ko{ㅅ}\\ \ko{ㅅ} drops\\ \ko{나아요}}}
    child {node[t] {\ko{르}\\ \ko{르} $\to$ \ko{ㄹㄹ}\\ \ko{몰라요}}}
    child {node[t] {\ko{ㅡ}\\ \ko{ㅡ} drops\\ \ko{아파요}}}
    child {node[t] {\ko{ㅎ}\\ \ko{ㅎ} drops,\\ vowel $\to$ \ko{ㅐ}\\ \ko{어때요}}}
    child {node[t] {\ko{ㄹ}\\ drops before\\ \ko{ㄴ ㅂ ㅅ}\\ \ko{사세요}}};
  \node[font=\scriptsize, text width=14cm, align=center] at (0,-3.2) {Only \emph{some} verbs ending in \ko{ㅂ ㄷ ㅅ ㅎ} are irregular (memorize them); \emph{all} verbs ending in \ko{르}, \ko{ㅡ} and \ko{ㄹ} follow their pattern (with very few exceptions).};
\end{tikzpicture}
\end{center}

\section{The patterns one by one}

\subsection{\ko{ㅂ} irregular}

\ko{ㅂ} becomes \ko{우} before a vowel; \ko{우} + \ko{어} contracts to \ko{워}. Almost
every \emph{adjective} ending in \ko{ㅂ} is irregular (\ko{덥다, 춥다, 쉽다, 어렵다,
맵다, 가깝다, 무겁다, 가볍다, 귀엽다, 반갑다}). The exceptions \ko{돕다} (help) and
\ko{곱다} (fine, lovely) give \ko{도와요}, \ko{고와요}.

\textbf{Regular}: \ko{입다} (wear) \ko{입어요}, \ko{잡다} (catch) \ko{잡아요},
\ko{좁다} (narrow) \ko{좁아요}, \ko{씹다} (chew) \ko{씹어요}.

\subsection{\ko{ㄷ} irregular}

\ko{ㄷ} becomes \ko{ㄹ} before a vowel. Common irregulars: \ko{듣다} (listen), \ko{걷다}
(walk), \ko{묻다} (ask), \ko{싣다} (load).

\textbf{Regular}: \ko{받다} (receive) \ko{받아요}, \ko{닫다} (close) \ko{닫아요},
\ko{믿다} (believe) \ko{믿어요}, \ko{묻다} (bury) \ko{묻어요} --- note the two different
verbs \ko{묻다}.

\subsection{\ko{ㅅ} irregular}

\ko{ㅅ} drops before a vowel, but the syllable does \emph{not} contract afterwards:
\ko{낫다} $\to$ \ko{나아요} (not \ko{*나요}). Common irregulars: \ko{낫다} (get
better), \ko{짓다} (build; make up a name), \ko{붓다} (swell; pour), \ko{젓다} (stir).

\textbf{Regular}: \ko{웃다} (laugh) \ko{웃어요}, \ko{씻다} (wash) \ko{씻어요},
\ko{벗다} (take off clothes) \ko{벗어요}.

\subsection{\ko{르} irregular}

Before \ko{\til 아/어}, the \ko{ㅡ} of \ko{르} drops and an extra \ko{ㄹ} is added to the
preceding syllable, giving \ko{ㄹ라} or \ko{ㄹ러}. The choice of \ko{라/러} follows the
vowel before \ko{르}: \ko{모르다} $\to$ \ko{몰라요}, \ko{부르다} $\to$ \ko{불러요}.
Before \ko{\til (으)면} etc.\ nothing changes: \ko{모르면}.

\begin{center}
\begin{tikzpicture}[every node/.style={font=\small},
    b/.style={draw, minimum size=0.95cm}]
  \node[b, fill=gray!10] (a1) at (0,0) {\ko{모}};
  \node[b, fill=orange!20] (a2) at (0.95,0) {\ko{르}};
  \node[b, fill=blue!10] (a3) at (2.4,0) {\ko{아요}};
  \node at (1.7,0) {+};
  \node at (3.5,0) {$\Longrightarrow$};
  \node[b, fill=gray!10] (b1) at (4.6,0) {\ko{몰}};
  \node[b, fill=orange!20] (b2) at (5.55,0) {\ko{라}};
  \node[b, fill=blue!10] (b3) at (6.5,0) {\ko{요}};
  \draw[-{Stealth}, red!70!black] (0.95,-0.55) to[bend right=25] node[below, font=\scriptsize] {\ko{ㄹ} copies back onto \ko{모}} (4.6,-0.55);
  \node[font=\scriptsize, anchor=west] at (7.2,0) {\ko{몰라요} mollayo, ``I don't know''};
\end{tikzpicture}
\end{center}

Common \ko{르} verbs: \ko{모르다} (not know), \ko{부르다} (sing; call; be full),
\ko{빠르다} (be fast), \ko{다르다} (be different), \ko{고르다} (choose), \ko{자르다}
(cut), \ko{기르다} (raise, grow), \ko{마르다} (be dry, thin), \ko{누르다} (press).
Exceptions: \ko{따르다} (follow, pour) $\to$ \ko{따라요} and \ko{들르다} (drop by)
$\to$ \ko{들러요} follow the \ko{ㅡ} pattern instead.

\subsection{\ko{ㅡ} irregular (``\ko{ㅡ} deletion'')}

\ko{ㅡ} drops before \ko{\til 아/어}; the vowel of the syllable before it decides
\ko{아} or \ko{어}; one-syllable stems take \ko{어}: \ko{아프다} $\to$ \ko{아파요},
\ko{예쁘다} $\to$ \ko{예뻐요}, \ko{쓰다} $\to$ \ko{써요}, \ko{끄다} (turn off) $\to$
\ko{꺼요}. Before \ko{\til (으)면} the stem already ends in a vowel, so it is just
\ko{아프면}.

\subsection{\ko{ㅎ} irregular}

Colour adjectives and the demonstrative adjectives \ko{이렇다}, \ko{그렇다},
\ko{저렇다}, \ko{어떻다} lose their \ko{ㅎ}. Before \ko{\til 아/어} the stem vowel and
ending merge into \ko{ㅐ} (\ko{ㅒ} after \ko{ㅑ}); before \ko{\til (으)} the \ko{으}
disappears too.
\begin{itemize}
  \item \ko{어떻다} $\to$ \ko{어때요?} (how is it?), \ko{어땠어요?}, \ko{어떠면}
        (rare), \ko{어떤} (what kind of).
  \item \ko{그렇다} $\to$ \ko{그래요} (that's right; OK), \ko{그러면} / \ko{그럼}
        (then, in that case).
  \item \ko{빨갛다} $\to$ \ko{빨개요} (it's red); \ko{노랗다} $\to$ \ko{노래요};
        \ko{하얗다} $\to$ \ko{하얘요}.
\end{itemize}
\textbf{Regular}: \ko{좋다} \ko{좋아요}, \ko{놓다} (put) \ko{놓아요/놔요}, \ko{넣다}
(put in) \ko{넣어요}.

\subsection{\ko{ㄹ} stems (``\ko{ㄹ} deletion'')}

Strictly speaking \ko{ㄹ}-stem verbs are regular before \ko{\til 아/어}, but their
\ko{ㄹ} drops before an ending starting with \ko{ㄴ}, \ko{ㅂ} or \ko{ㅅ} (mnemonic:
``\ko{ㄹ} hates \ko{ㄴ ㅂ ㅅ}''), and they never take the extra \ko{으}.
\begin{itemize}
  \item \ko{살다}: \ko{살아요} (regular), \ko{사세요} (\ko{ㅅ}), \ko{삽니다}
        (\ko{ㅂ}), \ko{사는 곳} (\ko{ㄴ}), \ko{살면}, \ko{살 거예요} (the future
        \ko{ㄹ} merges with the stem's \ko{ㄹ}).
  \item Other common ones: \ko{알다} (know), \ko{만들다} (make), \ko{놀다} (play),
        \ko{길다} (be long), \ko{멀다} (be far), \ko{열다} (open), \ko{팔다} (sell),
        \ko{돌다} (turn), \ko{걸다} (hang; dial).
\end{itemize}

\section{Summary table}

\begin{table}[H]
\centering
\small
\begin{tabular}{llllllll}
\toprule
Class & Verb & \ko{\til 아/어요} & \ko{\til 았/었어요} & \ko{\til (으)면} & \ko{\til (으)세요} & \ko{\til 고} & Meaning \\
\midrule
\ko{ㅂ} & \ko{덥다} & \ko{더워요} & \ko{더웠어요} & \ko{더우면} & \ko{더우세요} & \ko{덥고} & hot \\
\ko{ㅂ} & \ko{돕다} & \ko{도와요} & \ko{도왔어요} & \ko{도우면} & \ko{도우세요} & \ko{돕고} & help \\
\ko{ㄷ} & \ko{듣다} & \ko{들어요} & \ko{들었어요} & \ko{들으면} & \ko{들으세요} & \ko{듣고} & listen \\
\ko{ㄷ} & \ko{걷다} & \ko{걸어요} & \ko{걸었어요} & \ko{걸으면} & \ko{걸으세요} & \ko{걷고} & walk \\
\ko{ㅅ} & \ko{낫다} & \ko{나아요} & \ko{나았어요} & \ko{나으면} & \ko{나으세요} & \ko{낫고} & recover \\
\ko{ㅅ} & \ko{짓다} & \ko{지어요} & \ko{지었어요} & \ko{지으면} & \ko{지으세요} & \ko{짓고} & build \\
\ko{르} & \ko{모르다} & \ko{몰라요} & \ko{몰랐어요} & \ko{모르면} & \ko{모르세요} & \ko{모르고} & not know \\
\ko{르} & \ko{부르다} & \ko{불러요} & \ko{불렀어요} & \ko{부르면} & \ko{부르세요} & \ko{부르고} & sing, call \\
\ko{ㅡ} & \ko{아프다} & \ko{아파요} & \ko{아팠어요} & \ko{아프면} & \ko{아프세요} & \ko{아프고} & be sick \\
\ko{ㅡ} & \ko{크다} & \ko{커요} & \ko{컸어요} & \ko{크면} & \ko{크세요} & \ko{크고} & be big \\
\ko{ㅎ} & \ko{빨갛다} & \ko{빨개요} & \ko{빨갰어요} & \ko{빨가면} & --- & \ko{빨갛고} & be red \\
\ko{ㄹ} & \ko{살다} & \ko{살아요} & \ko{살았어요} & \ko{살면} & \ko{사세요} & \ko{살고} & live \\
\midrule
reg. & \ko{입다} & \ko{입어요} & \ko{입었어요} & \ko{입으면} & \ko{입으세요} & \ko{입고} & wear \\
reg. & \ko{받다} & \ko{받아요} & \ko{받았어요} & \ko{받으면} & \ko{받으세요} & \ko{받고} & receive \\
reg. & \ko{웃다} & \ko{웃어요} & \ko{웃었어요} & \ko{웃으면} & \ko{웃으세요} & \ko{웃고} & laugh \\
\bottomrule
\end{tabular}
\end{table}

\begin{important}[Exam strategy]
On the final, irregular verbs appeared in fill-in-the-blank conjugation questions. The
instructor's advice: (1) look at the last letter of the stem; (2) look at the first
letter of the ending --- if it is a consonant, do nothing (except for \ko{ㄹ}-stems);
(3) if it is a vowel, apply the pattern, then apply the usual vowel harmony and
contraction rules. The most-tested words were \ko{덥다, 춥다, 듣다, 걷다, 모르다,
부르다, 아프다, 바쁘다, 낫다, 어떻다, 살다, 만들다}.
\end{important}

\begin{example}[Irregulars in context]
One sentence for each class:
\begin{itemize}
  \item \ko{김치찌개가 너무 매워요.} Gimchijjigaega neomu maewoyo. (The kimchi stew is
        too spicy.) --- \ko{ㅂ}.
  \item \ko{매일 30분씩 걸어요.} Maeil samsip bunssik georeoyo. (I walk thirty minutes
        every day.) --- \ko{ㄷ}.
  \item \ko{감기가 다 나았어요.} Gamgiga da naasseoyo. (My cold is all better.) ---
        \ko{ㅅ}.
  \item \ko{그 사람 이름을 몰라요.} Geu saram ireumeul mollayo. (I don't know that
        person's name.) --- \ko{르}.
  \item \ko{요즘 너무 바빠요.} Yojeum neomu bappayo. (I'm so busy these days.) ---
        \ko{ㅡ}.
  \item \ko{한국어 수업은 어땠어요?} Hangugeo sueobeun eottaesseoyo? (How was Korean
        class?) --- \ko{ㅎ}.
  \item \ko{어디 사세요?} Eodi saseyo? (Where do you live?) --- \ko{ㄹ}.
\end{itemize}
\end{example}

\section{Course review: clause connectors and endings}

The second half of the course added a large set of endings. This table, which I made for
the final exam, collects them all in one place.

\begin{table}[H]
\centering
\begin{tabular}{llll}
\toprule
Pattern & Meaning & Example & Lesson \\
\midrule
\ko{\til 았/었어요} & past & \ko{봤어요} & 1 \\
\ko{\til 고 싶다} & want to & \ko{보고 싶어요} & 2 \\
\ko{\til 아/어 주세요} & please do for me & \ko{보여 주세요} & 2 \\
\ko{\til 고} & and; and then & \ko{맑고 따뜻해요} & 3 \\
\ko{\til 지만} & but & \ko{덥지만 좋아요} & 3 \\
\ko{\til 네요} & (I notice) & \ko{춥네요} & 3 \\
\ko{\til (으)시\til} & subject honorific & \ko{가세요, 읽으셨어요} & 4 \\
\ko{\til 는 것} & doing (nominalizer) & \ko{요리하는 것} & 5 \\
\ko{\til (으)ㄹ 수 있다/없다} & can / cannot & \ko{갈 수 있어요} & 5 \\
\ko{\til (으)러 가다} & go in order to & \ko{먹으러 가요} & 5 \\
\ko{\til (으)ㄹ 거예요} & will, going to; probably & \ko{갈 거예요} & 6 \\
\ko{\til (으)ㄹ까요?} & shall we/I? & \ko{갈까요?} & 6 \\
\ko{\til (으)ㅂ시다} & let's & \ko{갑시다} & 6 \\
\ko{\til (으)ㄹ게요} & I'll (promise) & \ko{갈게요} & 6 \\
\ko{\til 고 있다} & be \ldots ing & \ko{가고 있어요} & 7 \\
\ko{\til 아서/어서} & because; and then & \ko{늦어서} & 7 \\
\ko{\til (으)니까} & since (+ commands OK) & \ko{비가 오니까} & 7 \\
\ko{\til (으)면} & if, when & \ko{아프면} & 9 \\
\ko{\til 아/어야 하다} & must & \ko{가야 해요} & 9 \\
\ko{\til 아/어도 되다} & may & \ko{가도 돼요} & 9 \\
\ko{\til (으)면 안 되다} & must not & \ko{가면 안 돼요} & 9 \\
\ko{\til 지 마세요} & don't & \ko{가지 마세요} & 9 \\
\bottomrule
\end{tabular}
\end{table}

\begin{aside}
\textbf{Food culture review.} Many of the term's culture notes came back to food, which
in Korea is tied up with relationships and etiquette: wait for the eldest person to
start eating before you pick up your spoon; don't stick chopsticks upright in rice (it
recalls offerings to the dead); pour drinks for others rather than yourself, holding the
bottle with two hands for an older person, and turn your head away from elders when you
drink. The everyday greeting \ko{밥 먹었어요?} (have you eaten?) is less a question than
a way of saying ``how are you'', and side dishes (\ko{반찬} banchan) are shared and
refilled for free. Our end-of-term class party at a Korean restaurant in Waterloo was a
chance to practise all of it.
\end{aside}

\section*{Review summary}
\begin{itemize}
  \item Irregular changes happen only before vowel-initial endings (except
        \ko{ㄹ}-deletion before \ko{ㄴ ㅂ ㅅ}).
  \item \ko{ㅂ} $\to$ \ko{우} (\ko{더워요}); \ko{ㄷ} $\to$ \ko{ㄹ} (\ko{들어요});
        \ko{ㅅ} drops, no contraction (\ko{나아요}); \ko{르} $\to$ \ko{ㄹㄹ}
        (\ko{몰라요}); \ko{ㅡ} drops (\ko{아파요}); \ko{ㅎ} drops with \ko{ㅐ}
        (\ko{어때요}); \ko{ㄹ} drops before \ko{ㄴ ㅂ ㅅ} (\ko{사세요}).
  \item Learn the regular look-alikes: \ko{입다, 잡다, 받다, 닫다, 웃다, 씻다, 좋다}.
\end{itemize}
